\documentclass{article}
\usepackage[utf8]{inputenc}
\usepackage[T1]{fontenc}
\usepackage[a4paper,total={6in,8in}]{geometry}
\usepackage{amsmath,amssymb,amsfonts}
\usepackage{graphicx}
\usepackage{hyperref}
\usepackage{enumitem}
\usepackage{listings}
\usepackage{xcolor}
\DeclareMathOperator{\Var}{Var}
\DeclareMathOperator{\Cov}{Cov}
\providecommand{\tightlist}{\setlength{\itemsep}{0pt}\setlength{\parskip}{0pt}}
\lstset{basicstyle=\ttfamily\scriptsize,breaklines=true,columns=fullflexible,
  keepspaces=true,showstringspaces=false,frame=single}
\hypersetup{colorlinks=true,urlcolor=blue,linkcolor=blue}
\setlength{\emergencystretch}{2em}

\title{Math 126 Problem 17: Insights from the Edge-incidence Matrix}
\author{} % Put your name here.
\date{Spring 2027}
\begin{document}
\maketitle

\hypertarget{submission-instructions}{%
\subsection*{Submission Instructions}\label{submission-instructions}}

Please upload your work to Gradescope by \textbf{11:59pm on Monday, April 12, 2027 (Pacific time).}

\begin{itemize}
\item
  You are encouraged to discuss ideas and work with your classmates. However, you \emph{must acknowledge} your collaborators at the top of each solution on which you collaborated with others and you \emph{must write} your solutions independently.
\item
  Your solutions to theory questions must be written legibly, or typeset in LaTeX or markdown.
\item
  I recommend that you write your solutions in a Python notebook using Google Colab.
\item
  You should submit your solutions as a \emph{single PDF} via the assignment on Gradescope.
\end{itemize}

\textbf{Grading:} The point of the problem set is for \emph{you} to learn. To this end, I hope to disincentivize the use of LLMs by \emph{not} grading your work for correctness. Instead, you will grade your own work by comparing it to my solutions. This self-grade is due the Friday \emph{after} the problem set is due, also on Gradescope.

\noindent Original notebook: \href{https://colab.research.google.com/drive/16aOkT_-N2JFFow9si7DhIcgfnASXi8wt?usp=sharing}{Google Colab}.

% Source notebook cell 2
\hypertarget{insights-from-the-edge-incidence-matrix}{%
\subsection*{Insights from the Edge-incidence Matrix}\label{insights-from-the-edge-incidence-matrix}}

Let the number of edges \(m\) be greater than the number of nodes \(n\). The edge-incidence matrix \(\mathbf{B} \in \mathbb{R}^{m \times n}\) is an incredibly powerful tool for analyzing graph topology. In the notes, we saw that the Laplacian can be written as \(\mathbf{L} = \mathbf{B}^\top \mathbf{B}\). In this problem, we will look at the solutions to \(\mathbf{B}\mathbf{x} = \mathbf{0}\) and \(\mathbf{B}^\top\mathbf{y} = \mathbf{0}\) to uncover fundamental structural properties of the graph \(G=(V,E)\).

\hypertarget{part-a-node-vectors-and-connected-components}{%
\subsection*{Part A: Node Vectors and Connected Components}\label{part-a-node-vectors-and-connected-components}}

Consider a node vector \(\mathbf{x} \in \mathbb{R}^n\) such that \(\mathbf{B} \mathbf{x} = \mathbf{0}\). Observe that \(\mathbf{B}\mathbf{x} = \mathbf{0}\) if and only if \(\mathbf{L}\mathbf{x} = \mathbf{0}\).

Use the definition of \(\mathbf{B}\) to show that if \(\mathbf{B}\mathbf{x} = \mathbf{0}\), then \(x_i = x_j\) for every edge \((i,j) \in E\). Use this observation to formally prove that if the graph has \(k\) connected components, the multiplicity of the eigenvalue \(0\) of \(\mathbf{L}\) is exactly \(k\).

\hypertarget{part-b-edge-vectors-and-cycles}{%
\subsection*{Part B: Edge Vectors and Cycles}\label{part-b-edge-vectors-and-cycles}}

Now consider an edge vector \(\mathbf{y} \in \mathbb{R}^m\) such that \(\mathbf{B}^\top \mathbf{y} = \mathbf{0}\). We can think of each entry \(y_e\) as a ``flow'' assigned to the edge \(e\), respecting the arbitrary direction assigned to \(e\) by the definition of \(\mathbf{B}\). First, interpret the equation \(\mathbf{B}^\top \mathbf{y} = \mathbf{0}\) at a specific node \(v \in V\). What physical conservation law does this equation represent for the ``flows'' passing through \(v\)?

Next, recall from linear algebra that the number of independent solutions to a system of equations is equal to the number of variables minus the rank of the matrix. By Part A, we know that the rank of \(\mathbf{B}\) is \(n-k\) where \(k\) is the number of connected components. Using this fact, write down the number of independent solutions to \(\mathbf{B}^\top \mathbf{y} = \mathbf{0}\) in terms of \(m\), \(n\), and \(k\). Based on this formula and your physical interpretation of the flows, prove that there exists a non-zero solution \(\mathbf{y}\) if and only if the graph \(G\) contains at least one cycle.

\end{document}
